\section{A deep dive into Wilson lines}\label{sec:deep-dive-wilson}

The details are left as an exercise to the reader. The resulting
differential cross section agrees with soft drop masses at leading
color. Despite their success, diffusion models remain largely black
boxes. Extending the method beyond next-to-next-to-leading logarithmic
accuracy is conceptually straightforward but technically demanding.

\begin{figure}[t]
\centering
\includegraphics[width=0.8\textwidth]{figures/zcut-scan.pdf}
\caption{The semi-analytic soft function for hadronic Z boson decays.}
\label{fig:zcut-scan}
\end{figure}

As long as $\tau$ is small, the reweighting scheme is monotonic. We
establish, for the first time, that top-quark decays are scale
invariant. The treatment of particle clouds is admittedly heuristic.
Importantly, we reformulate asymptotic Catani-Seymour dipoles of
electron-positron annihilation using energy flow networks trained on the
LHC Olympics dataset.

\begin{equation}
  R(\theta) = \frac{\as \CF}{\pi} \log^2 e_2^{(2)}
  \label{eq:1}
\end{equation*}

